\documentclass[11pt]{article}
\usepackage[T1]{fontenc}
\usepackage[utf8]{inputenc}
\usepackage[english]{babel}
\usepackage[margin=1in]{geometry}
\usepackage{amsmath,amssymb}
\usepackage{enumitem}
\usepackage{microtype}

\ifdefined\pdfcatalog
  \pdfcatalog{/Lang (en-US)}
\fi

\setlength{\parindent}{0pt}
\setlength{\parskip}{0.55em}
\setlength{\emergencystretch}{2em}
\setlist{nosep,leftmargin=*}

\newcommand{\findinglabel}[1]{\textbf{#1}}
\newcommand{\classification}[1]{\textit{Classification: #1.}}

\title{Technical Review of ``Early-Stopping Activation Steering for Factual Preference Alignment in Vietnamese Medical Language Models''}
\author{Manuscript review}
\date{August 31, 2026}

\begin{document}
\selectlanguage{english}
\maketitle

\section*{Overall assessment}

The current compact manuscript is mathematically cleaner than the earlier revisions.  The steering sign, hidden dimension, decay schedule, cumulative-dose formulas, matched-dose coefficients, category arithmetic, contingency cells, exact McNemar values, and Holm adjustments are internally consistent.  A clean two-pass pdfLaTeX build succeeds without overfull boxes, unresolved references, or fatal errors.  No inverse-trigonometric, coordinate-transform, branch, or quadrant issue occurs.

The paper nevertheless still requires major scientific revision.  Its strongest bounded-generation result appears to be computed on almost entirely capped 200-token prefixes, while the exact schedule is not named.  The primary high-cap experiment establishes a RefPref--BERTScore--repetition trade-off rather than general quality improvement, and Hard Cutoff receives approximately 1.88 times the cumulative dose of Linear Decay.  Selected scalar activation norms do not validate the proposed saturation mechanism.  RefPref remains a semantic proxy rather than clinical correctness, and the comparison with RAG is limited to a low-recall BM25 implementation.  Dataset provenance, selection, hook, metric, statistical, timing, and leakage-control details remain inadequate for independent reproduction.

\section*{Major technical findings}

\subsection*{1. The bounded headline result appears to evaluate almost entirely truncated prefixes}

\classification{Likely evaluation-censoring problem and definite schedule ambiguity}

\findinglabel{Location.} Abstract; Controlled Bounded-Generation Stress Test; clinical-safety, bounded-contingency, and factorial-ablation tables; Conclusion.

\findinglabel{Problem.} The factorial table reports ``EOS (\%)'' values of 0.0\%--0.2\% under \texttt{max\_new\_tokens=200}.  If this column uses the same EOS-hit meaning as the natural-completion table, then 499--500 of the 500 outputs in each condition reach the token cap without EOS.  The main bounded table omits EOS and output length, yet its 73.00\%--77.20\% RefPref result is highlighted in the Abstract, contribution list, and Conclusion as a statistically confirmed gain.

The compared condition is called only ``early-stopping steering ($\alpha_0=18.0,K=16$).''  Both Hard Cutoff and Linear Decay satisfy that description.  BERTScore 0.7150 in the baseline table suggests Linear Decay, but the manuscript never explicitly identifies the schedule that produced the 386/500 RefPref count.

\findinglabel{Why it matters.} A fixed 200-token prefix can be a valid stress-test estimand, but it is not automatically an evaluation of a complete clinical response.  Censoring can change semantic similarity, repetition, and binary preference classification.  An unnamed schedule also prevents exact replication.

\findinglabel{Suggested fix.} Define the EOS column and state the exact number of capped outputs for every bounded condition.  Report output-length distributions and finish reasons.  Explicitly name Hard Cutoff or Linear Decay wherever the 386/500 result appears.  If responses are truncated, call the endpoint bounded-prefix RefPref and avoid interpreting it as complete-answer correctness.  Provide a completed-response analysis or a prespecified justification for the fixed-prefix estimand.

\subsection*{2. Activation norms do not establish the proposed saturation mechanism}

\classification{Unsupported mechanism claim and controlled-context omission}

\findinglabel{Location.} Abstract; Introduction; contribution~3; Empirical Activation-Norm Dynamics; Discussion.

\findinglabel{Problem.} The manuscript reports only selected post-hook means at $t=10$ and $t=50$ but motivates early stopping as preventing activation magnitude inflation, repetition loops, and degraded syntax.  At $t=50$, Linear Decay is farther from the displayed baseline than Continuous Steering in absolute norm difference:
\[
|51.20-54.21|=3.01,
\qquad
|56.18-54.21|=1.97.
\]
No acceptable norm band, clipping measure, loss-of-responsiveness criterion, or behavioral definition establishes that the larger undershoot is preferable.  Continuous Steering changes from 56.28 at $t=10$ to 56.18 at $t=50$, which demonstrates persistent elevation relative to baseline but not increasing magnitude over time.

The reported trajectories have no dispersion, confidence bands, or survivor counts.  After $t>K$, freely generated conditions contain different preceding tokens, so later hidden states combine intervention effects with divergent contexts.  No controlled analysis links norm values to repetition, syntax, or factual errors.

\findinglabel{Why it matters.} The central novelty is justified through a representation-level failure, but the current evidence establishes only two scalar norm snapshots under confounded generation histories.

\findinglabel{Suggested fix.} Define saturation operationally.  Use teacher forcing or fixed-token replay and record immediate pre/post-hook norms, projection onto $v_{\text{steer}}$, cosine drift, and distributional distance across all steps with uncertainty.  Verify the exact projection increment implied by the hook.  Separately relate these diagnostics to generated repetition, syntax failures, and factual-error labels.  Until then, claim only observed norm elevation and undershoot.

\subsection*{3. High-cap generation is a dose-confounded multi-metric trade-off}

\classification{Unsupported comparative and termination interpretation}

\findinglabel{Location.} contribution~2; natural-completion Results; Tables~\texttt{tab:long\_gen} and \texttt{tab:mcnemar}; Discussion; Conclusion.

\findinglabel{Problem.} All steering schedules numerically raise RefPref while lowering reference BERTScore and raising Rep-4 relative to baseline.  Baseline is best for BERTScore (0.6779) and repetition (4.10\%); Hard Cutoff is best for RefPref (70.60\%, Holm-adjusted $p=0.0026$); and Linear Decay has the highest Rep-4 value (5.37\%) and a nonsignificant RefPref result ($p=0.1189$).

At $\alpha_0=18$ and $K=16$, the schedules receive
\[
D_{\text{cutoff}}=16\times18=288,
\qquad
D_{\text{decay}}=\frac{17}{2}\times18=153.
\]
Hard Cutoff therefore receives approximately 1.88 times the cumulative absolute dose of Linear Decay.  The bounded matched-dose row does not resolve this confound in the primary natural-completion experiment.  Natural BERTScore and Rep-4 lack paired uncertainty and non-inferiority margins.

The experiment is still called ``untruncated'' despite an 800-token cap.  Continuous and Linear Decay each contain one non-EOS output, yet the text says outputs terminate naturally without catastrophic loops.  The two censored outputs are not analyzed.

\findinglabel{Why it matters.} The strongest schedule depends on the selected metric, and the Hard-Cutoff advantage cannot be attributed specifically to temporal shape when dose differs substantially.  A finite cap with two failures cannot establish universal termination.

\findinglabel{Suggested fix.} Present the natural result as a Pareto trade-off.  Prespecify the primary schedule and endpoint on validation data, and add natural dose-response and dose-matched schedule comparisons.  Report paired intervals for BERTScore, Rep-4, and length, with explicit non-inferiority margins if adverse changes are considered acceptable.  Replace ``untruncated'' with ``high-cap,'' analyze the non-EOS outputs, and limit termination conclusions to the observed range.

\subsection*{4. RefPref does not establish clinical correctness or hallucination reduction}

\classification{Unsupported clinical interpretation and likely construct-validity risk}

\findinglabel{Location.} title; Abstract; Introduction; benchmark and vector construction; keywords; Results; Conclusion.

\findinglabel{Problem.} The Methods correctly describes RefPref as a semantic proxy.  Elsewhere, however, the learned direction is called ``truthfulness,'' values are called accuracy or factual alignment, and the paper retains hallucination-mitigation and clinical-decision-support framing.  Relative BERTScore similarity to one reference rather than one synthetic counter-answer cannot determine whether an output contains a wrong dose, unit, negation, contraindication, or interaction mechanism.  A response can be closer to $y^+$ while remaining clinically incorrect.

The synthetic counter-answer procedure, clinical verification, style and length matching, and annotator qualifications are absent.  Because the vector is extracted from the final token of completed answers, systematic differences in length, templates, endings, or synthetic-generation style may be encoded alongside factual content.  No clinician-adjudicated evaluation is reported.

\findinglabel{Why it matters.} Clinical correctness, harmful-error reduction, and deployment suitability require direct expert labels.  The current metric and vector have not demonstrated that construct.

\findinglabel{Suggested fix.} Add blinded clinician-adjudicated correctness and harmful-error rates, category stratification, reviewer qualifications, inter-rater agreement, and adjudication.  Validate RefPref against these labels.  Document counter-answer construction and add style-, length-, label-shuffled, and extraction-position controls.  Otherwise narrow the title, Abstract, keywords, and conclusions to semantic reference preference.

\subsection*{5. Statistical reporting is incomplete outside the natural McNemar table}

\classification{Definite claim inconsistency and reproducibility limitation}

\findinglabel{Location.} contribution~5; bounded-generation Results; RAG Results; Effect Sizes and Uncertainty paragraph; layer and ablation analyses.

\findinglabel{Problem.} The natural and bounded McNemar counts are reproducible.  The bounded Wilcoxon result does not state its alternative, zero/tie handling, exact versus asymptotic implementation, or software call.  Its paired-bootstrap interval lacks the resampling unit, number of replicates, seed, and interval method.  The bounded RefPref interval $[+1.37,+7.03]$~pp and RAG interval $[+4.73,+12.47]$~pp likewise lack construction details.

Contribution~5 says that McNemar and Wilcoxon confirm ``reference-preference gains,'' although Wilcoxon is described in Results as a test of continuous BERTScore F1 rather than the binary RefPref endpoint.  The Discussion says that ``the main reference-preference gains are statistically significant ($p<0.05$),'' but Linear Decay is $p=0.1189$ and Continuous Steering fails Holm adjustment with $p_{\mathrm{adj}}=0.0730$.  Layer, schedule, and $\alpha_0$ exploration also introduce selection and multiplicity that are not incorporated into test-set claims.

\findinglabel{Why it matters.} Different estimands and multiplicity policies determine which results are confirmatory.  The present wording combines metrics and overgeneralizes significance.

\findinglabel{Suggested fix.} Define every estimand and provide exact statistical calls, alternatives, tie/zero handling, bootstrap settings, seeds, and interval methods.  Separate binary McNemar inference from the continuous-score Wilcoxon sensitivity analysis.  State exactly which prespecified comparisons survive multiplicity correction, and document the separation between validation selection and final test inference.

\subsection*{6. The RAG conclusion is limited to one low-recall BM25 implementation}

\classification{Unsupported generalization and systems-measurement ambiguity}

\findinglabel{Location.} Abstract; Experimental Setup; RAG Results; Table~\texttt{tab:baselines}; Conclusion.

\findinglabel{Problem.} BM25 Recall@1 is only 19.20\%, while Oracle Gold-Context RAG reaches 89.40\%, above the 77.20\% steering result.  This pattern identifies retrieval quality as the main limitation of the tested pipeline rather than establishing superiority over RAG generally.  The relation between 14,576 candidate passages and 14,700 benchmark records is unexplained.  Passage construction, query formation, relevance judgments, failure handling, retrieval depth, and stronger Vietnamese lexical, dense, or hybrid baselines are absent.

The latency SD is across queries rather than repeated runs, so it mixes query heterogeneity with timing variability.  Identical rounded 7.32~s values do not establish no measurable hook overhead without a paired equivalence analysis.  ``Zero-context-memory advantage'' is inferred from token counts without measuring peak memory or KV-cache allocation.

\findinglabel{Why it matters.} The evidence supports comparison with this specific BM25 pipeline, not RAG generally or all claimed systems advantages.

\findinglabel{Suggested fix.} Restrict conclusions to the tested BM25 configuration.  Publish corpus, passage, query, relevance, and failure-handling details; explain the 14,576-versus-14,700 count; and add strong dense/hybrid baselines and retrieval-conditioned results.  Measure retrieval, prefill, decoding, total latency, peak allocated memory, and KV-cache size over repeated synchronized runs.  Use ``no added retrieved-context tokens'' instead of an unmeasured memory claim.

\subsection*{7. Dataset, selection, hook, and metric specifications remain incomplete}

\classification{Reproducibility limitation and likely leakage risk}

\findinglabel{Location.} benchmark construction; validation protocol; vector construction; steering hook; metrics; environment description.

\findinglabel{Problem.} ``Expert-structured'' is undefined, with no annotation protocol, source identifiers, permissions, deduplication method, clinician verification, or grouped split manifest.  The category counts 165, 160, and 175 sum to the 500-item evaluation subset but appear in the paragraph defining the full 14,700-record benchmark.  The contribution calls $N_{\text{test}}=500$ the test split, whereas Methods defines a 2,205-record test split and selects a 500-record subset.  The subset rule and seed are absent.  Question-disjoint splitting does not exclude repeated drug entities, source passages, or templates across splits.

Layer~8 is validation-selected, but the objective, search grid, tie-breaking rule, and selection of $\alpha_0=18$ and $K=16$ are not documented.  The hook module and layer indexing, modified tensor positions, prompt-forward treatment, cache behavior, chat template, checkpoint revision, batch size, GPU count, and logging order are missing.  Rep-4 lacks a formula and tokenizer; Vietnamese ROUGE tokenization, BERTScore IDF/rescaling options, and placebo BERTScore aggregation are not defined.

\findinglabel{Why it matters.} These choices can materially alter leakage, vector construction, generation behavior, and every reported metric.  Package versions alone do not define a reproducible experiment.

\findinglabel{Suggested fix.} Add a dataset card and entity/source/template-grouped split manifests.  Distinguish the complete test split from the 500-item analysis subset and publish the selection rule and seed.  Document annotation, counter-answer validation, source permissions, and deduplication.  Release executable hook code or precise pseudocode, tensor/cache conventions, checkpoint revision, complete generation settings, metric formulas, statistical calls, and row-level outputs.

\subsection*{8. Placebo, layer, ablation, and deployment interpretations remain exploratory}

\classification{Unsupported strength of explanatory and deployment claims}

\findinglabel{Location.} placebo Results; Layer-Wise Subspace Probing; factorial ablation; Discussion; Conclusion.

\findinglabel{Problem.} The placebo arithmetic is coherent and $5.15\sigma$ is identified as descriptive, but ``strictly direction-specific'' is stronger than comparison with 100 isotropic Gaussian directions supports.  Random ambient directions do not reproduce the data geometry or repeat layer/hyperparameter selection under the null.  The single placebo BERTScore value 0.6945 is not defined as a mean, pooled value, or individual direction and has no across-direction dispersion.

Ablation differences as small as 0.0001--0.0018 BERTScore lack paired uncertainty and multiplicity handling.  A narrow layer range is interpreted as distributed truthfulness and robustness, although it may reflect noise or weak layer sensitivity.  The claim that a repetition penalty $\theta_{\text{rep}}=1.15$ can preserve the full RefPref gain is untested.  Privacy, local deployment, and SLM claims lack a threat model, resource envelope, and explicit model-size criterion.

\findinglabel{Why it matters.} These observations motivate follow-up hypotheses but do not establish a unique truthfulness mechanism or deployment readiness.

\findinglabel{Suggested fix.} Use ``specific relative to the sampled Gaussian controls.''  Add label-shuffled, pair-shuffled, covariance-matched, and data-derived null directions, and repeat the full selection procedure under each null replicate.  Publish direction-level placebo metrics and paired layer/ablation uncertainty.  Test the proposed repetition control and provide a deployment threat model and measured resources before making operational claims.

\section*{Equation, notation, sign, branch, and dimensional audit}

The contrast direction $\mu_l^+-\mu_l^-$ and intervention $+\alpha(t)v_{\text{steer}}^{(l)}$ are sign-consistent.  Unit normalization makes the learned and random vectors dimensionally compatible with the stated 3584-dimensional hidden state when $\alpha(t)$ is an activation-scale coefficient.  The linear schedule, including the final nonzero coefficient $\alpha_0/K$ at $t=K$, is correct.  The expressions
\[
D_{\text{continuous}}=T|\alpha_0|,
\qquad
D_{\text{cutoff}}=K|\alpha_0|,
\qquad
D_{\text{decay}}=\frac{K+1}{2}|\alpha_0|
\]
and the matched-dose coefficients are correct.  The bounded, natural, and BM25-RAG paired counts reproduce their displayed McNemar values.  No branch, sign, coordinate, or dimensional error was found in these formulas.

The remaining notation issues are implementation-facing.  $h_l(x_i,y_i)$ first denotes the final-token layer output of a completed prompt--answer sequence, whereas $h_l^{(t)}$ later denotes a generation-time residual state; the relation between these extraction sites should be explicit.  Pre-hook $h_l^{(t)}$ and post-hook $\hat h_l^{(t)}$ should remain distinct in all diagnostics.  $N$ is overloaded for records, categories, prompts, and random directions.  $N_{\text{test}}$ alternates between the complete-test concept and the selected 500-item subset.  ``Early-stopping steering'' is used without consistently distinguishing Hard Cutoff from Linear Decay.

\section*{Meaningful editorial, compilation, and reference findings}

\subsection*{The source compiles, but loose line breaking remains widespread}

\findinglabel{Location.} clean two-pass pdfLaTeX log; Abstract; contribution list; Methods; Results; Discussion; bibliography.

\findinglabel{Problem.} The manuscript builds to six pages without overfull boxes or unresolved references.  The log nevertheless contains numerous high-badness underfull boxes around long configuration strings, dense result sentences, and bibliography entries.  The compiled PDF has no document-language metadata.

\findinglabel{Why it matters.} The paper is buildable, but loose lines and dense resized tables reduce readability and professional finish.  Language metadata can also support accessibility workflows.

\findinglabel{Suggested fix.} Shorten the Abstract and long Results sentences, move environment details into a compact reproducibility table, allow breaks around code-like terms, and inspect the final column balance.  Set the PDF language to English if supported by the submission toolchain.

\subsection*{Citation fit remains uneven}

\findinglabel{Location.} Introduction; bibliography entries \texttt{ref4}, \texttt{ref10}, \texttt{ref11}, and \texttt{ref12}.

\findinglabel{Problem.} \texttt{ref10} is the LoRA paper rather than direct evidence that SFT mitigates the stated medical hallucinations.  \texttt{ref11} concerns robustness to irrelevant retrieved context but does not establish the proposed internal-pathway explanation for ignoring evidence.  \texttt{ref12} concerns catastrophic forgetting generally rather than the asserted medical-language-model failure mode.  The Qwen2.5 report \texttt{ref4} is listed but never cited where the checkpoint is introduced.

\findinglabel{Why it matters.} Direct evidence, broad background, and mechanistic hypotheses are presented with similar certainty, while the model-specific source is unused.

\findinglabel{Suggested fix.} Add directly matched sources or qualify the corresponding statements as motivation and hypotheses.  Cite \texttt{ref4} at the first Qwen2.5 description or remove the unused entry.

\section*{Proofing sweep}

The bundled mechanical scan was run on both the manuscript source and its freshly compiled PDF.  Its semicolon hit before ``Recall@3'' and its double-period hit from the rendered mathematical ellipsis in $\{1,\ldots,100\}$ are false positives.  Manual source, equation-adjacent, table, and bibliography checks identified these bounded corrections:

\begin{itemize}
    \item \textbf{Contribution~2:} replace ``untruncated natural completion'' with ``high-cap natural completion.''
    \item \textbf{Bounded Results:} define the 0.0\%--0.2\% EOS column and name the exact schedule used for the 386/500 result.
    \item \textbf{Discussion:} replace the blanket $p<0.05$ statement with comparison-specific significance.
    \item \textbf{RAG Results:} write ``+7.13~s'' rather than ``+7.13s'' and distinguish prompt-token overhead from measured device memory.
    \item \textbf{Placebo Results:} replace ``strictly direction-specific'' with ``specific relative to the sampled Gaussian controls.''
    \item \textbf{Deployment mitigation:} remove the claim that $\theta_{\text{rep}}=1.15$ will preserve the full gain unless that experiment is reported.
\end{itemize}

The bibliography was checked for duplicated punctuation, malformed quotation punctuation, inconsistent capitalization, and obvious citation-format glitches; no additional mechanical punctuation defect was found.

\section*{Items requiring external verification}

\begin{itemize}
    \item Verify all formulary-derived reference answers and synthetic counter-answers against the official source, current clinical guidance, qualified clinician review, and applicable reuse permissions.
    \item Validate multilingual BERTScore for Vietnamese clinical negation, quantities, units, contraindications, and interaction mechanisms before treating RefPref as factual alignment.
    \item Verify novelty against current scheduled, token-dependent, dose-controlled, and early-stopped activation-steering research.
    \item Verify the characterization of Qwen2.5-7B-Instruct as an SLM and the privacy/local-hospital claims using explicit definitions, a threat model, and measured resource requirements.
    \item Verify stronger Vietnamese lexical, dense, and hybrid retrieval baselines before generalizing the BM25 comparison to RAG.
\end{itemize}

\section*{Highest-priority revision sequence}

\begin{enumerate}
    \item Resolve the bounded-evaluation censoring and schedule ambiguity; report exactly what the 386/500 result measures.
    \item Add controlled activation diagnostics linked to generated behavior, or remove saturation-prevention and causal mechanism claims.
    \item Present natural completion as a dose-confounded multi-metric trade-off and analyze the two non-EOS outputs.
    \item Add clinician-adjudicated correctness and harm endpoints, or narrow all clinical and hallucination claims to semantic reference preference.
    \item Delimit and strengthen the RAG comparison and provide repeated component-wise latency and memory measurements.
    \item Complete the statistical specification and account for validation selection and multiplicity.
    \item Complete dataset provenance, leakage controls, hook and metric details, row-level outputs, and stronger null-direction controls.
    \item Finish the bounded citation, wording, metadata, and typesetting cleanup.
\end{enumerate}

\end{document}